\section{结果讨论}

\subsection{静止系密度与 $Z$ 因子}



一个重要的对象是静止系密度 ${\cal M} (0,z_3^2)$。
它由 $P=0$ 的数据产生。
其虚部的结果如预期与零相容。
其实部示于图 \ref{M0}，是 $z_3$ 的偶函数，
并且在较小和中等 $|z_3|$ 范围内，其随 $|z_3|$ 的衰减中可以清楚地看到线性成分。
事实上，线性指数因子 $Z(z_3^2) \sim  e^{-c |z_3|/a}$ 的出现
正是直线规范链接产生非微扰效应的表现。


\begin{figure}[t]
  \centerline{\includegraphics[width=3.7in]{images/M0NormK2.pdf}}
    \caption{静止系密度 ${\cal M} (0,z_3^2)$ 的实部
    \label{M0}}
    \end{figure}


\subsection{约化 Ioffe 时间分布}

图 \ref{realz} 中我们画出了比值实部
${\cal M} (Pz_3, z_3^2)/{\cal M} (0,z_3^2)$ 作为 $z_3$ 的函数的结果，
取六个固定的动量 $P$ 值。
可以看到所有曲线都具有类高斯的形状。
因此，$Z(z_3^2)$ 链接重正化因子已在比值中被消去，符合预期。



此外，各条曲线彼此相似，{仅随 $P$ 增大宽度递减}。
在图 \ref{realc} 中，我们对同一组数据改用 $\nu=Pz_3$ 作横轴。
可以看到，现在数据实际上落在同一条曲线上。
虚部的情形也类似。

这一现象对应于软 TMD ${\cal F} (x, k_\perp^2)$ 中
$x$ 依赖与 $k_\perp$ 依赖的因子化，正如前面几节所讨论的。


\begin{figure}[t]
    \centerline{\includegraphics[width=4.0in]{images/ReMz.pdf} }
    \caption{约化分布 ${\mathfrak  M} (Pz_3,z_3^2)$ 的实部作为 $z_3$ 的函数。
    这里 $P= 2 \pi p/L$。
    \label{realz}}
    \end{figure}


        \begin{figure}[t]
    \centerline{\includegraphics[width=3.8in]{images/RealC.pdf}}
    \caption{${\mathfrak  M} (\nu, z_3^2)$ 的实部作为 $\nu =Pz_3$ 的函数，
    并与式 (\relax {\ref{MC}})、(\relax {\ref{qV}}) 给出的曲线比较。
        \label{realc}}
    \end{figure}

      \begin{figure}[t]
    \centerline{\includegraphics[width=3.8in]{images/pheno.pdf} }
    \caption{由 \mbox{式 (\relax {\ref{qV}})} 给出的价分布 $q_v(x) $
    与 $Q^2=1$ GeV$^2$ 处的 NNLO 全局拟合
    NNPDF31\_nnlo\_pch\_as\_0118\_mc\_164 \protect{\cite{Ball:2017nwa}}、
    MSTW2008nnlo68cl\_nf4 \protect{\cite{Martin:2009iq}} 以及 NLO 全局拟合 CJ15nlo
    \protect{\cite{Accardi:2016qay}} 的比较；
    三者均通过 LHAPD6 库 \protect{\cite{Buckley:2014ana}} 抽取。
    全局拟合周围的能带表示其实验与系统不确定度。
        \label{MSTW}}
    \end{figure}


    \subsection{夸克—反夸克分解}

Ioffe 时间分布的实部由如下余弦 Fourier 变换得到：
  \begin{align}
{\mathfrak  M}_R (\nu) \equiv   & \int_0^1 dx \,
\cos (\nu x)  \, q_v  (x)
\label{MC}
     \end{align}
其中函数 $q_v(x)$ 是夸克与反夸克分布之差
$q_v(x) = q(x) -\bar q(x)$。
在我们的情形，$q$ 即 $u-d$，
$\bar q = \bar u - \bar d$。
$u-\bar u$ 的 $x$ 积分等于质子内 $u$ 夸克的数目，即 2；
而 $d- \bar d$ 的 $x$ 积分等于 1。
于是 $q_v (x)$ 的 $x$ 积分应等于 1。



我们发现，若选取函数
     \begin{align}
    q_v(x) =  \frac{315}{32} \sqrt{x}  (1-x)^{3}\,  ,
    \label{qV}
    \end{align}
（其 $x$ 积分归一化为 1），则实部数据可以得到很好的描述。
具体做法是：对归一化的 $x^a(1-x)^b$ 型函数作余弦 Fourier 变换
${\cal M} (\nu;a,b)$，
通过拟合数据确定参数 $a,b$。
数据与基于式 (\relax {\ref{MC}})、(\relax {\ref{qV}}) 曲线的比较见图 \ref{realc}。

虽然拟合使用了全部数据点，但后者明显由
Re ${\mathfrak  M} (\nu, z_3^2)$ 取值较小的那些点主导。
当 $\nu <10$ 时，位于曲线上方的数据点对应
$z_3=3a$ 到 $5a$ 的取值。稍后我们将看到，
它们反映的是微扰演化：
\mbox{Re ${\mathfrak  M} (\nu, z_3)$} 随 $z_3$ 减小而增大。
在此意义上，整体曲线 (\ref{qV}) 对应“低归一化点”处的 PDF，
即微扰演化停止的区域。

一般而言，更恰当的做法是把 Re ${\mathfrak  M} (\nu, z_3^2)$
作为两个变量 $\nu$ 与 $z_3$ 的函数来拟合，
尽管 $z_3$ 依赖相当微弱，仅对少数几个点可辨。
由于我们只把 Re ${\mathfrak  M} (\nu, z_3)$ 作为单一变量 $\nu$ 的函数进行拟合，
确实存在一些明显偏离曲线的点，但我们认为
把小 $z_3$ 点的演化 $z_3$ 依赖折算成
图 \mbox{\ref{realc}} 中曲线的误差带并无意义。
在第 5.4 节中，我们把数据点演化到共同参考标度 $z_0=2a$，
并展示由此所得结果的误差带。

我们清楚本格点设定相当粗糙（淬火近似、非常大的 π 介子质量），
因此不打算将结果与实验数据做细致比较。
不过我们认为某种程度的比较作为示例仍相当有用。

于是，我们把 $q_v (x) $ 与价分布之差 \mbox{$u_v(x) -d_v(x)$} 的三个全局拟合进行比较，
见图 \ref{MSTW}。
这些全局拟合曲线对应 $\mu=1$ GeV 标度，
而我们的“低归一化点”曲线对应 $\mu \lesssim 0.3$ GeV。
尽管如此，可以看到我们的曲线并不太远：
与 NNPDF31 \cite{Ball:2017nwa} NNLO 拟合直到 $x=0.1$ 都相去不远，
与 MSTW \cite{Martin:2009iq} NNLO 拟合直到 $x=0.05$ 也较为接近。
我们还给出了 NLO 拟合 CJ15 \cite{Accardi:2016qay}。




由于每条曲线下的面积都等于 1，我们的曲线通过在 $x>0.1$ 区间超出 NNLO 曲线，
补偿了它在 $x<0.1$ 区域的强烈不足。
换言之，若我们的曲线在 $x<0.1$ 区域描述得更好，
则它在 $x>0.1$ 区域必然会更小。


      \begin{figure}[t]
    \centerline{\includegraphics[width=3.9in]{images/ImC.pdf} }
    \caption{${\mathfrak  M} (\nu, z_3^2)$ 的虚部
    与基于 $\bar q (x)=0$ 的曲线 $  {\mathfrak  M}_I^v (\nu)$ 的比较。
        \label{imc}}
    \end{figure}


      \begin{figure}[t]
    \centerline{ \includegraphics[width=3.9in]{images/ImCA.pdf}}
    \caption{${\mathfrak  M} (\nu, z_3^2)$ 的虚部
    与基于式 (\relax {\ref{qA}}) 所给 $\bar q(x)$ 的曲线的比较。
        \label{imca}}
    \end{figure}



正弦 Fourier 变换
      \begin{align}
  {\mathfrak  M}_I (\nu)  \equiv
& \int_0^1 dx \,
\sin (\nu x) \, q_+  (x)
     \end{align}
由函数 $q_+ (x)= q(x) + \bar q (x)$ 构成，
它也可写成 \mbox{$q_+ (x)= q_v(x) + 2 \bar q (x)$}。
如果忽略反夸克贡献而取 \mbox{$q_+  (x) = q_v(x)$}，
便得到图 \ref{imc} 所示的曲线（记作 $  {\mathfrak  M}_I^v (\nu)$）。
若采用非零的反夸克贡献
        \begin{align}
    \bar q(x) =\bar u (x) - \bar d(x) =  0.07 \left [ 20  \,  x (1-x)^{3} \right ] \,  ,
    \label{qA}
    \end{align}
则与数据的符合显著改善，见图 \ref{imca}。
该函数是这样得到的：用 $A x^a (1-x)^b$ 函数的正弦 Fourier 变换
去拟合差值 ${\rm Im} \ {\mathfrak M} (\nu, z_3^2) - {\mathfrak M}_I^v(\nu)$ 的数据。
此结果对应于
          \begin{align}
 \int_0^1 dx \, [\bar u (x) - \bar d(x)] =  0.07 \ .
     \end{align}
定义在 $-1\leq x \leq 1$ 区间上的组合分布
         \begin{align}
q(x)& = u(x)-d(x)    \nn & =  [q_v (x) +\bar q (x)  ]\, \theta(x>0)
-
\bar q (-x)  \, \theta(x<0)
\label{q}
     \end{align}
示于图 \ref{umind}。


          \begin{figure}[t]
    \centerline{\includegraphics[width=4.0in]{images/umindkey.pdf} }
    \caption{由式 (\relax {\ref{q}}) 定义的总体分布 $q(x)$。
        \label{umind}}
    \end{figure}


        \subsection{演化}


数据与一条 \mbox{$z_3$ 无关}的曲线总体上相符固然令人满意，
但仍容易注意到数据中的残余 \mbox{$z_3$ 依赖}。
当某个特定 $\nu$ 对应多个不同 $z_3$ 取值的数据点时，
这一点尤其明显。
检验这种依赖是否对应微扰演化是很有意思的。




首先，实部的演化应导致其随 $z_3^2$ 增大而{\it 减小}。
另一方面，如第 3 节末尾所指出的，
只要 \mbox{$\nu \lesssim 5.5$}，
函数 ${{\rm Im} \, \cal M} ( \nu, z_3^2)$ 就随 $z_3^2$ {\it 增大}而增大。
我们的数据遵循这些模式。

如前所述，演化对应于 Ioffe 时间分布在小的 $z_3^2$ 处的 $\ln z_3^2$ 奇异性。
因此一个自然的想法是检验：对应于小 $z'_3$ 与 $z_3$ 的数据
是否可以通过
  \begin{align}
{\mathfrak M} (\nu, {z'}_3^2) {=}
{\mathfrak M} (\nu, z_3^2) \
 -\frac23  \frac{\alpha_s}{\pi}   \ln ({z'_3}^2/z_3^2) B \otimes {\mathfrak M}\,  (\nu, z_3^2)
 \label{Prop}
 \end{align}
对某一 $\alpha_s$ 取值相互关联起来。
此处 $B$ 是演化核 (\ref{Bu})。
在我们的情形，
    \begin{align}
B \otimes {\mathfrak M}\,  (\nu)   & =
\int_0^1  du \,  \frac{1+u^2}{1- u} \,    [{\mathfrak M} ( \nu)- {\mathfrak  M} (u \nu)  ]
 .
\label{plusM}
 \end{align}





更具体地说，我们把点 $z'_3$ 固定在 \mbox{$z_0=2a$}——
在领头对数精度下它对应
$\overline {\rm MS}$ 方案标度 \mbox{$\mu_0 = 1$  GeV }——
并用不同的 $\alpha_s$ 取值，
从 ${\mathfrak M}\,  (\nu, z_3^2)$ 的数据点构造函数
   \begin{align}
\widetilde{\mathfrak M} (\nu, z_0^2) {\equiv }
{\mathfrak M} (\nu, z_3^2) \
 -\frac23  \frac{\alpha_s}{\pi}   \ln (z_0^2/z_3^2) B \otimes {\mathfrak M}\,  (\nu, z_3^2)
 \label{EvoM}
 \end{align}

由于预期微扰演化适用于小的 $z_3$，
我们在分析中纳入了 $z_3$ 直至 4 个格距的数据，
对应的能量标度为 \mbox{$\mu = 2,     1,     0.7$  和     0.5  GeV}。


                \begin{figure}[t]
    \centerline{ \includegraphics[width=3.9in]{images/EvoRe0.pdf}}
    \caption{$z_3/a =1, 2, 3, 4$ 时 ${\mathfrak  M} (\nu, z_3^2)$ 的实部。
        \label{ER0}}
    \end{figure}

      \begin{figure}[t]
    \centerline{ \includegraphics[width=3.8in]{images/EvoReEv.pdf}}
    \caption{实部演化后的数据点。
        \label{ERE}}
    \end{figure}



实部的这些数据点示于 \mbox{图 \ref{ER0}}。
可以看到数据点有明显弥散。
取 $\alpha_s/\pi =0.1$，我们计算出对应函数
$\widetilde{\mathfrak M} (\nu, z_0^2) $ 的“演化后”数据点。
结果示于图 \ref{ERE}。演化后的数据点现在非常接近一条普适曲线。



    \begin{figure}[t]
    \centerline{ \includegraphics[width=3.8in]{images/EvoIm0.pdf}}
    \caption{$z_3/a =1, 2, 3, 4$ 时 ${\mathfrak  M} (\nu, z_3^2)$ 的虚部。
        \label{EI0}}
    \end{figure}

           \begin{figure}[t]
    \centerline{ \includegraphics[width=3.8in]{images/EvoImEv.pdf}}
    \caption{虚部演化后的数据点。
        \label{EIE}}
    \end{figure}


        \begin{figure}[t]
    \centerline{ \includegraphics[width=4.0in]{images/Revol10.pdf}}
    \caption{$z_3 \leq 10a$ 的 Re ${\mathfrak M}\,  (\nu, z_3^2)$ 数据点按文中所述方法演化到 $z_0=2a$。
        \label{evpoi}}
    \end{figure}


        \begin{figure}[t]
    \includegraphics[width=4.2in]{images/evol4.pdf}
    \caption{
    由图 \ref{evpoi} 所示演化后数据构造的 $u_v(x) - d_v(x)$ 曲线，
    视为对应 \mbox{$\mu^2 =1$  GeV$^2$} 标度；
    再将其演化到全局拟合的参考点 $\mu^2 =4$  GeV$^2$。
        \label{evolved}}
    \end{figure}



在图 \ref{EI0} 中我们给出虚部的初始数据点；
用同一 $\alpha_s/\pi =0.1$ 数值构造的演化后数据点示于图 \ref{EIE}。
它们同样接近一条普适曲线。
这一分析表明：固定 $\nu$ 下 ${\mathfrak M}\,  (\nu, z_3^2)$ 的残余 $z_3^2$ 依赖
与小 $z_3^2$ 处预期的对数演化相容。
显然这是我们计算的一个重要特征，需要进一步研究，
因为它将在从这类格点计算可靠抽取重正化 PDF 的过程中发挥关键作用。

若采用更小的格距，微扰演化的适用范围可以在更宽的 $\nu$ 区域得到辩护。
虽然我们的数据延伸到 $\sim 1$ fm 的较大间隔，
但我们发现把它们当作展示微扰演化所生趋势的一个例子是有益的。

为此，我们对 $z_3 \leq 6a$ 的数据点应用了领头对数公式 (\ref{EvoM})
（取 $z_0=2a$、$\alpha_s/\pi=0.1$）。
假设演化在 $z_3 \gtrsim 6a$ 处停止（如我们的数据所示），
则 \mbox{$7a \leq z_3 \leq 10a$} 的数据点用 $z_3=6a$ 代入式 (\ref{EvoM}) 演化到 $z_0$。
如此演化后的数据点示于图 \ref{evpoi}。


用归一化 $N(a,b) x^a(1-x)^b$ 型函数的余弦 Fourier 变换
${\cal M} (\nu;a,b)$ 拟合演化后的点，
我们发现取 $a = 0.36(6)$、$b=3.95(22)$ 即可描述它们。
把 $z_0=2a$ 视为 $\overline{\rm MS}$ 标度 $\mu=1$ GeV，
便可进一步把该曲线演化到全局拟合的标准参考标度 $\mu^2 = 4$ GeV$^2$，
见图 \ref{evolved}。
与图 \ref{MSTW} 比较，可见微扰演化移动了我们的曲线，
使其更靠近全局拟合。
